% Abstract for the PacificVis 2027 conference-track abstract deadline (Nov 2, 2026 AoE), System Paper category.
% PCS takes the title and abstract as plain text; a plain-text copy is at the end of this file.
% Build separately from main.tex: pdflatex -output-directory=build_abstract abstract_submission.tex
\documentclass[11pt]{article}
\usepackage[margin=1in]{geometry}
\usepackage{times}

\title{Read, Trace, Edit: Probing the Residual Stream with Two Lenses}
\author{}
\date{}

\begin{document}
\maketitle

\begin{abstract}
Lenses that decode a language model's hidden states into words are a popular way to look inside it, but a heatmap
of decoded words shows what a vector \emph{reads as}, not whether the model uses it. We present an interactive viewer
that puts reading, tracing and editing on one layers-by-tokens grid of the residual stream. Every cell can be decoded
with two lenses, the logit lens and the Jacobian lens (J-lens), traced back to the tokens it drew from, and edited,
with the layers above rerun on the spot. Three usage scenarios on Gemma-3-4B, each backed by a sweep over hundreds of
prompts, show what this combination finds. First, the two lenses suit different edits. Pushing a word's logit-lens
direction into the last token makes the model say it from layer~6--9, while the J-lens direction needs layers 16--30.
Swapping one entity for another along J-lens directions redirects fact lookups (73\% of country--capital pairs, against
2\% for the logit lens) but fails on antonyms and translation (at most 7\%). The viewer shows why: a J-lens swap works
only when the J-lens can read the word in its own cell, which it does for names from the first layers and for words
with many senses only after the model has moved the word on. Second, in two-hop questions (the capital of the state
containing a given city), the state is readable at the city token from layer~7, and swapping it there redirects the
answer 92--98\% of the time; the same state, read at the answer position, is not used. Third, across 500
chain-of-thought questions in five task types, the final answer is often readable before the model writes any
reasoning (top five in 97\% of correct fact-retrieval answers), and the J-lens reads it at the same or an earlier layer
than the logit lens in every task type, much earlier for serial deduction (median layer 17 against 32). For ordering
puzzles, neither lens reads the answer before the reasoning (at most 11\%, no more than a control word), consistent
with the model working it out as it writes.
\end{abstract}

\end{document}

% ---- Plain text for the PCS form ----
% Lenses that decode a language model's hidden states into words are a popular way to look inside it, but a heatmap
% of decoded words shows what a vector reads as, not whether the model uses it. We present an interactive viewer that
% puts reading, tracing and editing on one layers-by-tokens grid of the residual stream. Every cell can be decoded
% with two lenses, the logit lens and the Jacobian lens (J-lens), traced back to the tokens it drew from, and edited,
% with the layers above rerun on the spot. Three usage scenarios on Gemma-3-4B, each backed by a sweep over hundreds
% of prompts, show what this combination finds. First, the two lenses suit different edits. Pushing a word's
% logit-lens direction into the last token makes the model say it from layer 6-9, while the J-lens direction needs
% layers 16-30. Swapping one entity for another along J-lens directions redirects fact lookups (73% of
% country-capital pairs, against 2% for the logit lens) but fails on antonyms and translation (at most 7%). The viewer
% shows why: a J-lens swap works only when the J-lens can read the word in its own cell, which it does for names from
% the first layers and for words with many senses only after the model has moved the word on. Second, in two-hop
% questions (the capital of the state containing a given city), the state is readable at the city token from layer 7,
% and swapping it there redirects the answer 92-98% of the time; the same state, read at the answer position, is not
% used. Third, across 500 chain-of-thought questions in five task types, the final answer is often readable before
% the model writes any reasoning (top five in 97% of correct fact-retrieval answers), and the J-lens reads it at the
% same or an earlier layer than the logit lens in every task type, much earlier for serial deduction (median layer 17
% against 32). For ordering puzzles, neither lens reads the answer before the reasoning (at most 11%, no more than a
% control word), consistent with the model working it out as it writes.
